\section{C++ ROS2 implementation update}

In the final implementation, all runtime ROS2 nodes were rewritten in \textbf{C++ using rclcpp}, while the
offline metrics aggregation and report-generation tools remained in \textbf{Python}. This split was intentional:
C++ was used for the online components involved in runtime sensing, coordination, logging, and benchmark
execution, whereas Python was retained for offline data analysis and automatic report generation.

\subsection{Runtime nodes implemented in C++}

\begin{table}[H]
\centering
\caption{Final implementation split between online C++ ROS2 nodes and offline Python tools.}
\small
\begin{tabularx}{\linewidth}{>{\bfseries}p{3.2cm}p{3.1cm}X>{\raggedright\arraybackslash}p{1.6cm}}
\toprule
Component & Role & CogAR interpretation & Language \\
\midrule
\texttt{braitenberg\_reflex} & immediate obstacle-driven turn-away response & bio-inspired behavior + computational + publish-subscribe & C++ \\
\texttt{subsumption\_mux} & fixed-priority command arbitration & subsumption / fixed-priority coordination & C++ \\
\texttt{benchmark\_logger} & odom/cmd/scan logging with service control & computational + request-process-reply & C++ \\
\texttt{benchmark\_runner} & scenario execution and metadata generation & plan/act interface for benchmark orchestration & C++ \\
\texttt{topological\_goal\_manager\_cli} & symbolic route computation over topological areas & hierarchical planning / symbolic GO-TO operator & C++ \\
\texttt{summarize\_results.py} & per-run metric extraction & offline computational analysis & Python \\
\texttt{aggregate\_results.py} & cross-run aggregation & experiment pipeline support & Python \\
\texttt{update\_report\_data.py} & generation of tables and figures for LaTeX & reproducibility support & Python \\
\bottomrule
\end{tabularx}
\end{table}

\subsection{Why this split is architecturally coherent}

This implementation split also reinforces the CogAR interpretation of the software. The online ROS2 nodes are
responsible for sensing, coordination, safety reaction, and benchmark execution, which makes them part of the
\textbf{robot cognitive architecture itself}. In contrast, the Python scripts operate after the run has completed and
therefore belong to the \textbf{experimental evaluation pipeline}, not to the runtime robot architecture.

\subsection{Build and reproducibility notes}

The project is meant to be built in a ROS2 Humble environment using \texttt{colcon build --symlink-install}. Dependencies
are resolved through \texttt{rosdep}, and the final tables/plots are generated from recorded CSV logs rather than inserted
manually. This keeps the report consistent with the benchmark evidence produced in Gazebo.

\subsection{Verification items to report}

The final report should explicitly mention which of the following checks were completed on the target machine:
\begin{itemize}
  \item successful package build with \texttt{colcon};
  \item correct registration of the package executables;
  \item visibility of the runtime nodes through \texttt{ros2 node list};
  \item correct advertisement of \texttt{/benchmark\_logger/enable};
  \item successful route generation by the topological planner CLI;
  \item creation of \texttt{run\_metadata.yaml} and benchmark CSV logs during an actual run.
\end{itemize}

\subsection{Limitations to acknowledge}

The C++ rewrite does not remove the remaining integration dependencies of the project. In particular, the final
Gazebo execution still depends on the lab-provided G1 package, the final locomotion command topic may require a
small adapter, and TEB should only be included in the benchmark if a Humble-compatible plugin is actually available
in the lab environment.
